\chapter{DGA / Ω 异常}
\chaptervoice{这一章，我们围绕“DGA / Ω 异常”做一件克制的事：把直觉压缩为状态、关系、时间尺度和可检验后果。它在全书中的任务，是说明持续主体的功能异常和恢复。}

\section*{理论来路与依赖接口}
本章的直接思想来源是原文关于控制工程、培育工程和心理工程的生命周期衍生讨论。我们采用原文较晚形成的定义来整理较早讨论，不把对话中的试探性比喻与最终模型并列成互相竞争的“定律”。已有理论锚点主要是系统辨识、记忆偏置、反馈失稳与慢变量保护（Kalman、Ashby 与混杂系统理论）。

本章使用上一章“自我与身份异常”已经得到的对象与边界；本章产出将供下一章“Agent 心理干预”使用。若后文术语偶尔出现，只作为路线提示，不参与本章推导。

\section*{从哪里出发}
我们已经拥有上一章给出的概念，现在只向前走一步。我会先把对象放进闭环，再讨论它的数学形式，最后检查工程边界。读者不必预先相信本章术语；只需暂时接受一个方法论约定：智能结构的意义来自它在现实反馈中的作用，而不是来自名称本身。

令系统在观察窗口 $[t,t+T]$ 内的状态轨迹为 $x_{[t,t+T]}$，环境扰动为 $d$，行动为 $u$。本章讨论的对象若确实有效，就应改变条件分布 $p(x_{t+T}\mid x_t,u,d)$，或改变达到同一结果所需的代价。这个起点把讨论锚定在系统辨识、记忆偏置、反馈失稳与慢变量保护上。
\section{长期方向结构失真}
\begin{narration}
现在把镜头移到“长期方向结构失真”。我们只沿本章已经取得的概念向前一步：先确定它在闭环里的角色，再检查它的尺度与失败边界。
\end{narration}

本节把“长期方向结构失真”落实为对未来可接受生命周期分布施加偏好的极慢生成吸引子。这个定义不是说“任何相似现象都算”，而是要求我们同时给出对象域、允许的干预以及观察窗口。对于主题“DGA / Ω 异常”而言，它承担的是持续主体的功能异常和恢复中的一个可辨识环节。

把这一关系写成最小数学骨架，可得

\begin{equation}p_\Omega(\xi)\propto e^{-V_\Omega(\xi)}.\end{equation}

式中的符号只承担结构角色，具体参数仍需由任务和身体数据辨识。我们首先检查可观测量，再检查控制量，最后检查比它更慢的约束。这样的顺序来自系统辨识、记忆偏置、反馈失稳与慢变量保护，但本书对长期方向结构失真的命名和组合仍是需要实验验证的建模提案。

这一小节的反例同样重要：把方向偏好写成未来对过去的逆因果会混淆生成先验与物理时间。因此评价它不能只看一次任务结果，而要比较干预前后的轨迹、恢复时间、维护成本和风险。如果改变该机制而所有允许观察下的分布都不变，那么它至少在当前实验族中是冗余描述。

\begin{thought}
对“长期方向结构失真”做一次消融：冻结它、删除它，或只允许它以相邻时间尺度交换残差。哪些现象立即消失，哪些只在长期出现？若答案无法区分三种干预，我们还没有找到正确的状态变量。
\end{thought}

最后，把结论交还现实：本节只建立功能层关系，不由此推出特定脑区实现、主观意识或宇宙目的。可测状态、干预后果和明确失败条件，才是从原文直觉走向可检验理论的桥。

\section{Goal Preference Drift}
\begin{narration}
现在把镜头移到“Goal Preference Drift”。我们只沿本章已经取得的概念向前一步：先确定它在闭环里的角色，再检查它的尺度与失败边界。
\end{narration}

本节把“Goal Preference Drift”落实为对未来可接受生命周期分布施加偏好的极慢生成吸引子。这个定义不是说“任何相似现象都算”，而是要求我们同时给出对象域、允许的干预以及观察窗口。对于主题“DGA / Ω 异常”而言，它承担的是持续主体的功能异常和恢复中的一个可辨识环节。

把这一关系写成最小数学骨架，可得

\begin{equation}p_\Omega(\xi)\propto e^{-V_\Omega(\xi)}.\end{equation}

式中的符号只承担结构角色，具体参数仍需由任务和身体数据辨识。我们首先检查可观测量，再检查控制量，最后检查比它更慢的约束。这样的顺序来自系统辨识、记忆偏置、反馈失稳与慢变量保护，但本书对Goal Preference Drift的命名和组合仍是需要实验验证的建模提案。

这一小节的反例同样重要：把方向偏好写成未来对过去的逆因果会混淆生成先验与物理时间。因此评价它不能只看一次任务结果，而要比较干预前后的轨迹、恢复时间、维护成本和风险。如果改变该机制而所有允许观察下的分布都不变，那么它至少在当前实验族中是冗余描述。

\begin{thought}
对“Goal Preference Drift”做一次消融：冻结它、删除它，或只允许它以相邻时间尺度交换残差。哪些现象立即消失，哪些只在长期出现？若答案无法区分三种干预，我们还没有找到正确的状态变量。
\end{thought}

最后，把结论交还现实：本节只建立功能层关系，不由此推出特定脑区实现、主观意识或宇宙目的。可测状态、干预后果和明确失败条件，才是从原文直觉走向可检验理论的桥。

\section{Self 与 DGA 不一致}
\begin{narration}
现在把镜头移到“Self 与 DGA 不一致”。我们只沿本章已经取得的概念向前一步：先确定它在闭环里的角色，再检查它的尺度与失败边界。
\end{narration}

本节把“Self 与 DGA 不一致”落实为跨经历保持身份连续、同时允许受约束慢漂移的不变量集合。这个定义不是说“任何相似现象都算”，而是要求我们同时给出对象域、允许的干预以及观察窗口。对于主题“DGA / Ω 异常”而言，它承担的是持续主体的功能异常和恢复中的一个可辨识环节。

把这一关系写成最小数学骨架，可得

\begin{equation}\dot\Psi=\varepsilon_\Psi\Pi_{\mathcal I}g(E,\Theta,\Psi).\end{equation}

式中的符号只承担结构角色，具体参数仍需由任务和身体数据辨识。我们首先检查可观测量，再检查控制量，最后检查比它更慢的约束。这样的顺序来自系统辨识、记忆偏置、反馈失稳与慢变量保护，但本书对Self 与 DGA 不一致的命名和组合仍是需要实验验证的建模提案。

这一小节的反例同样重要：更新过快导致身份碎裂，完全冻结则阻断必要成长。因此评价它不能只看一次任务结果，而要比较干预前后的轨迹、恢复时间、维护成本和风险。如果改变该机制而所有允许观察下的分布都不变，那么它至少在当前实验族中是冗余描述。

\begin{thought}
对“Self 与 DGA 不一致”做一次消融：冻结它、删除它，或只允许它以相邻时间尺度交换残差。哪些现象立即消失，哪些只在长期出现？若答案无法区分三种干预，我们还没有找到正确的状态变量。
\end{thought}

最后，把结论交还现实：本节只建立功能层关系，不由此推出特定脑区实现、主观意识或宇宙目的。可测状态、干预后果和明确失败条件，才是从原文直觉走向可检验理论的桥。

\section{过早固化}
\begin{narration}
现在把镜头移到“过早固化”。我们只沿本章已经取得的概念向前一步：先确定它在闭环里的角色，再检查它的尺度与失败边界。
\end{narration}

本节把“过早固化”落实为瞬时状态、带时间索引的经验轨迹与轨迹生成结构之间的双向相变。这个定义不是说“任何相似现象都算”，而是要求我们同时给出对象域、允许的干预以及观察窗口。对于主题“DGA / Ω 异常”而言，它承担的是持续主体的功能异常和恢复中的一个可辨识环节。

把这一关系写成最小数学骨架，可得

\begin{equation}h_t\rightarrow E_{[t-\Delta,t]}\xrightarrow{\mathsf C}\Theta.\end{equation}

式中的符号只承担结构角色，具体参数仍需由任务和身体数据辨识。我们首先检查可观测量，再检查控制量，最后检查比它更慢的约束。这样的顺序来自系统辨识、记忆偏置、反馈失稳与慢变量保护，但本书对过早固化的命名和组合仍是需要实验验证的建模提案。

这一小节的反例同样重要：把三相做成三个独立数据库会切断写入、召回和巩固的动力联系。因此评价它不能只看一次任务结果，而要比较干预前后的轨迹、恢复时间、维护成本和风险。如果改变该机制而所有允许观察下的分布都不变，那么它至少在当前实验族中是冗余描述。

\begin{thought}
对“过早固化”做一次消融：冻结它、删除它，或只允许它以相邻时间尺度交换残差。哪些现象立即消失，哪些只在长期出现？若答案无法区分三种干预，我们还没有找到正确的状态变量。
\end{thought}

最后，把结论交还现实：本节只建立功能层关系，不由此推出特定脑区实现、主观意识或宇宙目的。可测状态、干预后果和明确失败条件，才是从原文直觉走向可检验理论的桥。

\section{过度可塑}
\begin{narration}
现在把镜头移到“过度可塑”。我们只沿本章已经取得的概念向前一步：先确定它在闭环里的角色，再检查它的尺度与失败边界。
\end{narration}

本节把“过度可塑”落实为对未来可接受生命周期分布施加偏好的极慢生成吸引子。这个定义不是说“任何相似现象都算”，而是要求我们同时给出对象域、允许的干预以及观察窗口。对于主题“DGA / Ω 异常”而言，它承担的是持续主体的功能异常和恢复中的一个可辨识环节。

把这一关系写成最小数学骨架，可得

\begin{equation}p_\Omega(\xi)\propto e^{-V_\Omega(\xi)}.\end{equation}

式中的符号只承担结构角色，具体参数仍需由任务和身体数据辨识。我们首先检查可观测量，再检查控制量，最后检查比它更慢的约束。这样的顺序来自系统辨识、记忆偏置、反馈失稳与慢变量保护，但本书对过度可塑的命名和组合仍是需要实验验证的建模提案。

这一小节的反例同样重要：把方向偏好写成未来对过去的逆因果会混淆生成先验与物理时间。因此评价它不能只看一次任务结果，而要比较干预前后的轨迹、恢复时间、维护成本和风险。如果改变该机制而所有允许观察下的分布都不变，那么它至少在当前实验族中是冗余描述。

\begin{thought}
对“过度可塑”做一次消融：冻结它、删除它，或只允许它以相邻时间尺度交换残差。哪些现象立即消失，哪些只在长期出现？若答案无法区分三种干预，我们还没有找到正确的状态变量。
\end{thought}

最后，把结论交还现实：本节只建立功能层关系，不由此推出特定脑区实现、主观意识或宇宙目的。可测状态、干预后果和明确失败条件，才是从原文直觉走向可检验理论的桥。

\section{生命周期方向性崩溃}
\begin{narration}
现在把镜头移到“生命周期方向性崩溃”。我们只沿本章已经取得的概念向前一步：先确定它在闭环里的角色，再检查它的尺度与失败边界。
\end{narration}

本节把“生命周期方向性崩溃”落实为对未来可接受生命周期分布施加偏好的极慢生成吸引子。这个定义不是说“任何相似现象都算”，而是要求我们同时给出对象域、允许的干预以及观察窗口。对于主题“DGA / Ω 异常”而言，它承担的是持续主体的功能异常和恢复中的一个可辨识环节。

把这一关系写成最小数学骨架，可得

\begin{equation}p_\Omega(\xi)\propto e^{-V_\Omega(\xi)}.\end{equation}

式中的符号只承担结构角色，具体参数仍需由任务和身体数据辨识。我们首先检查可观测量，再检查控制量，最后检查比它更慢的约束。这样的顺序来自系统辨识、记忆偏置、反馈失稳与慢变量保护，但本书对生命周期方向性崩溃的命名和组合仍是需要实验验证的建模提案。

这一小节的反例同样重要：把方向偏好写成未来对过去的逆因果会混淆生成先验与物理时间。因此评价它不能只看一次任务结果，而要比较干预前后的轨迹、恢复时间、维护成本和风险。如果改变该机制而所有允许观察下的分布都不变，那么它至少在当前实验族中是冗余描述。

\begin{thought}
对“生命周期方向性崩溃”做一次消融：冻结它、删除它，或只允许它以相邻时间尺度交换残差。哪些现象立即消失，哪些只在长期出现？若答案无法区分三种干预，我们还没有找到正确的状态变量。
\end{thought}

最后，把结论交还现实：本节只建立功能层关系，不由此推出特定脑区实现、主观意识或宇宙目的。可测状态、干预后果和明确失败条件，才是从原文直觉走向可检验理论的桥。

\section*{本章收束：把名词还原为闭环}
现在回头看“DGA / Ω 异常”，最重要的不是记住缩写，而是说清本章对象怎样改变系统的状态、代价或可恢复性。下一章“Agent 心理干预”只使用这里已经给出的结果，不反过来为本章提供前提。

\begin{bookproposition}
若一个候选机制既不改变可观测轨迹分布，也不改变资源、风险或恢复代价，那么在当前观察尺度上，它不是一个可辨识的功能结构。
\end{bookproposition}
\begin{proofidea}
功能等价由可干预后果定义。若所有允许干预下的输出分布与代价均不变，则该机制相对于当前实验族不可辨识；要么扩大观察尺度，要么删除这一冗余描述。
\end{proofidea}

\begin{readercheck}
请尝试回答：本章对象的状态是什么？它的最快和最慢时间常数是什么？什么观测能证伪它？它失败时，残差应向哪一层升级？
\end{readercheck}
\cleardoublepage
